\chapter{电压与电压表}\label{ch:g8:voltage-voltmeter}

自第一块\omterm{def:g3:first-electric-circuit:battery}{电池}点亮第一只\omterm{def:g3:first-electric-circuit:bulb}{灯泡}起，本书中便有一个词承担着
沉重而含糊的任务：“推”。\omterm{def:g3:first-electric-circuit:battery}{电池}推得轻，市电推得凶，
两块首尾相接的\omterm{def:g3:first-electric-circuit:battery}{电池}合力而推。
“推”如今成为一个可\omterm{def:g6:measurement-in-science:measuring}{测量}的量，
拥有著名的\omterm{def:g6:measurement-in-science:measuring}{单位}，而它的仪器，与上一章的不同，
拒绝站在路中间。

\section{电压}

\begin{definition}[电压]\label{def:g8:voltage-voltmeter:voltage}
电路中两点之间的\emph{电压}\index{电压} $U$，
度量两点之间可用的电气“推力”——
即电流从一点被推向另一点的力度，
正如瀑布的落差度量它推动水轮的力度。
电压总是\emph{两点之间}的量：\omterm{def:g3:first-electric-circuit:battery}{电池}的两个\omterm{def:g3:first-electric-circuit:battery}{接线柱}之间，\omterm{def:g3:first-electric-circuit:bulb}{灯泡}的两端之间，
绝不只限于某一个点。
其\omterm{def:g6:measurement-in-science:measuring}{单位}是
\emph{伏特}\index{伏特}（\unit{V}），以纪念第一块\omterm{def:g3:first-electric-circuit:battery}{电池}的
发明者。
\end{definition}

\begin{example}[伏特家族相册]\label{ex:g8:voltage-voltmeter:values}
单节圆形\omterm{def:g3:first-electric-circuit:battery}{电池}：\qty{1.5}{V}。扁平的袖珍\omterm{def:g3:first-electric-circuit:battery}{电池}：
\qty{4.5}{V}——外包装内三节\omterm{def:g3:first-electric-circuit:battery}{电池}首尾相接。
长方形九伏\omterm{def:g3:first-electric-circuit:battery}{电池}，两个\omterm{def:g3:first-electric-circuit:battery}{接线柱}都在顶部：
\qty{9}{V}，内含六节隐藏的电芯。汽车\omterm{def:g3:first-electric-circuit:battery}{电池}：\qty{12}{V}。墙上的
插座：\qty{230}{V}——比一百五十多节小\omterm{def:g3:first-electric-circuit:battery}{电池}还要多的“推力”，
这正是“插座，绝不可碰”这句话的全部定量内容。
云与地之间的闪电：
高达数亿\omterm{def:g8:voltage-voltmeter:voltage}{伏特}，但转瞬即逝。
\end{example}

\begin{proposition}[无电流时的电压]\label{prop:g8:voltage-voltmeter:nocurrent}
\omterm{def:g8:voltage-voltmeter:voltage}{电压}与电流是两个不同的量，而且\omterm{def:g8:voltage-voltmeter:voltage}{电压}是先行的一方。
一节单独放在抽屉里、不接任何导线的\omterm{def:g3:first-electric-circuit:battery}{电池}，
不会驱动任何电流——但它两端之间
完整的\qty{1.5}{V}\omterm{def:g8:voltage-voltmeter:voltage}{电压}依然存在，可以\omterm{def:g6:measurement-in-science:measuring}{测量}，
就像水磨停工的那天，瀑布的落差依旧。
推力可以存在而无水流；水流却永远离不开推力。
\end{proposition}

\begin{definition}[电压表]\label{def:g8:voltage-voltmeter:voltmeter}
\emph{电压表}\index{电压表}\omterm{def:g6:measurement-in-science:measuring}{测量}两点之间的\omterm{def:g8:voltage-voltmeter:voltage}{电压}——
因此它应当\emph{跨接}在这两点之间：
与所研究的元件\emph{\omterm{def:g5:series-parallel-circuits:parallel}{并联}}，两根表笔分别接在两侧，
好比\omterm{def:g6:measurement-in-science:measuring}{测量}员把水准仪架在两个高度之间。
与\omterm{def:g8:intensity-ammeter:ammeter}{电流表}在各方面都截然相反，它几乎
\emph{完全}地阻碍电流：跨接在任何元件上，它几乎分不走
一点电流，因此接上它不会扰动任何东西，
也无需断开任何回路。
\end{definition}

\begin{omfigure}
\begin{tikzpicture}[scale=0.9]
  \draw (0,0) to[battery1, l=电池] (0,2.8)
    to[short] (2.0,2.8)
    to[lamp, l=灯] (4.6,2.8)
    to[short] (6.4,2.8)
    to[short] (6.4,0)
    to[short] (0,0);
  \draw (2.0,2.8) to[short] (2.0,4.6)
    to[voltmeter, l=$U$] (4.6,4.6)
    to[short] (4.6,2.8);
\end{tikzpicture}

{\small \omterm{def:g8:voltage-voltmeter:voltmeter}{电压表}采取\omterm{def:g6:measurement-in-science:measuring}{测量}员的姿态：跨接在
\omterm{def:g3:first-electric-circuit:bulb}{灯泡}两侧，每侧一根表笔，\omterm{def:g6:measurement-in-science:measuring}{测量}两端之间的“推力”——
绝不站在路中间。}
\end{omfigure}

\begin{method}[测量电压]\label{met:g8:voltage-voltmeter:measure}
\begin{enumerate}
  \item 保持电路完好——对\omterm{def:g8:voltage-voltmeter:voltmeter}{电压表}而言，
        无须断开任何部分；
  \item 将一根表笔接在元件的每一侧：
        \unit{V}表笔接在常规电流流入的一端，
        公共端接在电流流出的一端；
  \item 先用最大量程，再逐档调至合适量程；
  \item 读数，并写出\omterm{def:g6:measurement-in-science:measuring}{单位}。
\end{enumerate}
两表的“脾气”并列如下：\omterm{def:g8:intensity-ammeter:ammeter}{电流表}——断开电路，
站入其中，不设阻碍；
\omterm{def:g8:voltage-voltmeter:voltmeter}{电压表}——不断开电路，站在一旁，处处阻碍电流。混淆两者的
“脾气”，是实验室里代价最高的习惯。
\end{method}

\section{额定电压}

\begin{definition}[额定电压]\label{def:g8:voltage-voltmeter:nominal}
每个\omterm{def:g3:first-electric-circuit:bulb}{灯泡}、\omterm{ex:g6:circuits-and-safety:motor}{电动机}和电器都标有\emph{额定
电压}\index{额定电压}——即设备设计时所针对的\omterm{def:g8:voltage-voltmeter:voltage}{电压}，
印在底座或铭牌上：手电筒\omterm{def:g3:first-electric-circuit:bulb}{灯泡}“3.5 V”、汽车前灯“12 V”、
水壶“230 V”。以额定\omterm{def:g8:voltage-voltmeter:voltage}{电压}供电时，
设备按设计运行；供电远低于此，它便闹脾气——
\omterm{def:g3:first-electric-circuit:bulb}{灯泡}昏暗，\omterm{ex:g6:circuits-and-safety:motor}{电动机}无力；供电远高于此，它便烧毁——
亮得惊人，却只此一次。
\end{definition}

\begin{example}[灯泡与电池的匹配]\label{ex:g8:voltage-voltmeter:matching}
标有“3.5 V”的\omterm{def:g3:first-electric-circuit:bulb}{灯泡}接在\qty{4.5}{V}扁平\omterm{def:g3:first-electric-circuit:battery}{电池}上：略微过载，
灯光明亮，但寿命短暂。接在单节
\qty{1.5}{V}\omterm{def:g3:first-electric-circuit:battery}{电池}上：只有昏昏欲睡的橙色微光。接在\qty{9}{V}
\omterm{def:g3:first-electric-circuit:battery}{电池}上：刹那间耀眼，随后细丝断裂——
这是灯丝的讣告。标签即契约：遵守它，
\omterm{def:g3:first-electric-circuit:bulb}{灯泡}便能活出标称的寿命。
\end{example}

\begin{remark}[为何 230 V 会致命]\label{rem:g8:voltage-voltmeter:mains}
现在，那条古老的绝对规则有了它的数字。人体潮湿的
导电通路，接上\qty{4.5}{V}时，通过的电流
小得无法察觉。同一通路接上\qty{230}{V}——\omterm{def:g8:voltage-voltmeter:voltage}{电压}是原来的五十倍——
通过的电流也是五十倍，远超足以使心脏停跳的
百分之几\omterm{def:g8:intensity-ammeter:intensity}{安培}的电流。插座并非
“只是略强”于你的\omterm{def:g3:first-electric-circuit:battery}{电池}实验台：它比之高出
两个数量级，而规章的严厉程度，
不过是简单的算术而已。
\end{remark}

\section{初步测量电压}

\begin{example}[绕简单回路一周]\label{ex:g8:voltage-voltmeter:readings}
一节\omterm{def:g3:first-electric-circuit:battery}{电池}、一只\omterm{def:g3:first-electric-circuit:bulb}{灯泡}，手持电表。\omterm{def:g6:measurement-in-science:measuring}{测量}\omterm{def:g3:first-electric-circuit:battery}{电池}两端：
\qty{4.5}{V}。\omterm{def:g6:measurement-in-science:measuring}{测量}\omterm{def:g3:first-electric-circuit:bulb}{灯泡}两端：\qty{4.5}{V}——\omterm{def:g3:first-electric-circuit:bulb}{灯泡}
得到的正是\omterm{def:g3:first-electric-circuit:battery}{电池}提供的\omterm{def:g8:voltage-voltmeter:voltage}{电压}。\omterm{def:g6:measurement-in-science:measuring}{测量}一段
连接\emph{导线}两端：\qty{0}{V}，或电表所能示出的最接近的值——
示数。好导线如同平坦的道路，不耗费任何“推力”。
“推力”存在于\emph{用电器}两端，而不在导线上：
这一读数值得记住，待下一章的定律将它
正式确认。
\end{example}

\begin{example}[断开开关的“供认”]\label{ex:g8:voltage-voltmeter:openswitch}
\omterm{def:g3:first-electric-circuit:bulb}{灯泡}不亮，\omterm{ex:g3:first-electric-circuit:switch}{开关}断开——推力去了哪里？\omterm{def:g6:measurement-in-science:measuring}{测量}：
不亮的\omterm{def:g3:first-electric-circuit:bulb}{灯泡}两端为 \qty{0}{V}；两端的\emph{断开的
\omterm{ex:g3:first-electric-circuit:switch}{开关}}则为满额 \qty{4.5}{V}。\omterm{def:g3:first-electric-circuit:battery}{电池}的推力完整地
横跨缺口等待——回路中唯一一处
两点在电学上真正“相距甚远”。合上
\omterm{ex:g3:first-electric-circuit:switch}{开关}，两处读数互换。断开的\omterm{ex:g3:first-electric-circuit:switch}{开关}
正是\omterm{def:g8:voltage-voltmeter:voltage}{电压}等待的地方。
\end{example}

\section{练习}

\begin{exercise}[$\star$]\label{exo:g8:voltage-voltmeter:1}
\omterm{def:g8:voltage-voltmeter:voltage}{电压}\omterm{def:g6:measurement-in-science:measuring}{测量}的是什么？为什么它总是“两点之间”的量？
写出\omterm{def:g8:voltage-voltmeter:voltage}{电压}的\omterm{def:g6:measurement-in-science:measuring}{单位}，并说明水流比喻。
\end{exercise}

\begin{exercise}[$\star$]\label{exo:g8:voltage-voltmeter:2}
凭记忆写出：圆柱形\omterm{def:g3:first-electric-circuit:battery}{电池}、扁平\omterm{def:g3:first-electric-circuit:battery}{电池}、
方形\omterm{def:g3:first-electric-circuit:battery}{电池}、汽车\omterm{def:g3:first-electric-circuit:battery}{电池}和市电的\omterm{def:g8:voltage-voltmeter:voltage}{电压}。
\end{exercise}

\begin{exercise}[$\star$]\label{exo:g8:voltage-voltmeter:3}
一节\omterm{def:g3:first-electric-circuit:battery}{电池}未接任何用电器。有多大电流通过？
跨接其两端的\omterm{def:g8:voltage-voltmeter:voltmeter}{电压表}示数是多少？哪一条命题在此适用？
\end{exercise}

\begin{exercise}[$\star$]\label{exo:g8:voltage-voltmeter:4}
比较两种电表的用法：连接方式、是否“断路”、
对电流的阻碍作用。哪一种错误会借助表盘造成
\omterm{def:g7:short-circuits-safety:device}{短路}？
\end{exercise}

\begin{exercise}[$\star$]\label{exo:g8:voltage-voltmeter:5}
什么是\omterm{def:g8:voltage-voltmeter:nominal}{额定电压}？预测一个“6 V”\omterm{def:g3:first-electric-circuit:bulb}{灯泡}
在 \qty{1.5}{V}、\qty{6}{V}、\qty{12}{V} 下的表现。
\end{exercise}

\begin{exercise}[$\star$]\label{exo:g8:voltage-voltmeter:6}
在简单回路中，\omterm{def:g8:voltage-voltmeter:voltmeter}{电压表}跨接\omterm{def:g3:first-electric-circuit:battery}{电池}、\omterm{def:g3:first-electric-circuit:bulb}{灯泡}
以及一段导线两端时，分别读数多少？推力
存在于何处？
\end{exercise}

\begin{exercise}[$\star$]\label{exo:g8:voltage-voltmeter:7}
\omterm{ex:g3:first-electric-circuit:switch}{开关}断开时，不亮的\omterm{def:g3:first-electric-circuit:bulb}{灯泡}两端和断开的\omterm{ex:g3:first-electric-circuit:switch}{开关}两端各是多少？
合上\omterm{ex:g3:first-electric-circuit:switch}{开关}的瞬间，两处读数会怎样变化？
\end{exercise}

\begin{exercise}[$\star$]\label{exo:g8:voltage-voltmeter:8}
为何把\omterm{def:g8:voltage-voltmeter:voltmeter}{电压表}并接在正在工作的\omterm{def:g3:first-electric-circuit:bulb}{灯泡}两端，几乎不会扰动
电路？这只\omterm{def:g8:voltage-voltmeter:voltmeter}{电压表}拒绝
去做什么？
\end{exercise}

\begin{exercise}[$\star\star$]\label{exo:g8:voltage-voltmeter:9}
扁平的 \qty{4.5}{V} \omterm{def:g3:first-electric-circuit:battery}{电池}内藏三节单体\omterm{def:g3:first-electric-circuit:battery}{电池}。画出（或描述）
它们的排列方式，并计算若其中一节被误装反时的\omterm{def:g8:voltage-voltmeter:voltage}{电压}
——即多年前手电筒一章的内容，如今用数字来算。
（方向相反的推力相减。）
\end{exercise}

\begin{exercise}[$\star\star$]\label{exo:g8:voltage-voltmeter:10}
露营者的\omterm{def:g3:first-electric-circuit:bulb}{灯泡}额定 \qty{12}{V}，房车\omterm{def:g3:first-electric-circuit:battery}{电池}新充满时读数为
\qty{12.6}{V}，经过一个漫长的夜晚后为 \qty{11.2}{V}。
描述\omterm{def:g3:first-electric-circuit:bulb}{灯泡}的这一晚，并说明它为何没有烧坏。
\end{exercise}

\begin{exercise}[$\star\star$]\label{exo:g8:voltage-voltmeter:11}
误将\omterm{def:g8:voltage-voltmeter:voltmeter}{电压表}与\omterm{def:g3:first-electric-circuit:bulb}{灯泡}\emph{\omterm{def:g4:series-circuits:series}{串联}}接入（回路断开，电表接在缺口处）：
预测\omterm{def:g3:first-electric-circuit:bulb}{灯泡}的表现
与电表的示数——别忘了\omterm{def:g8:voltage-voltmeter:voltmeter}{电压表}所阻碍的是什么。
（\omterm{def:g3:first-electric-circuit:bulb}{灯泡}不亮与电表示数都很有启发意义，
值得细想。）
\end{exercise}

\begin{exercise}[$\star\star\star$]\label{exo:g8:voltage-voltmeter:12}
两只相同的\omterm{def:g3:first-electric-circuit:bulb}{灯泡}\omterm{def:g4:series-circuits:series}{串联}接在 \qty{4.5}{V} 的电源上：\omterm{def:g8:voltage-voltmeter:voltmeter}{电压表}
测得每只\omterm{def:g3:first-electric-circuit:bulb}{灯泡}两端约为 \qty{2.25}{V}。将这两只\omterm{def:g3:first-electric-circuit:bulb}{灯泡}
\omterm{def:g5:series-parallel-circuits:parallel}{并联}接在同一\omterm{def:g3:first-electric-circuit:battery}{电池}上：预测每只\omterm{def:g3:first-electric-circuit:bulb}{灯泡}的\omterm{def:g8:voltage-voltmeter:voltage}{电压}，并
用你多年来已知的亮度规律调和两种结果
（\omterm{def:g4:series-circuits:series}{串联}的一对较暗，\omterm{def:g5:series-parallel-circuits:parallel}{并联}的一对
全亮）。哪一章将使你的预测成为定律？
\end{exercise}

\section{问题：电池抽屉大检查}

\begin{problem}[{周末问题——电池抽屉的
大检查；一只电压表、一盒散装单体电池，以及家中的
小电器架}]\label{pb:g8:voltage-voltmeter:1}
家家都有这样一个抽屉：散落的单体\omterm{def:g3:first-electric-circuit:battery}{电池}、来历不明的\omterm{def:g3:first-electric-circuit:battery}{电池}，
以及“需要换新”的小电器。检查日——\omterm{def:g8:voltage-voltmeter:voltmeter}{电压表}
在手。

\medskip\noindent\textbf{第一部分——单体\omterm{def:g3:first-electric-circuit:battery}{电池}分类。}
\begin{enumerate}
  \item 描述用\omterm{def:g8:voltage-voltmeter:voltmeter}{电压表}检查一节
        散放的 \qty{1.5}{V} 单体\omterm{def:g3:first-electric-circuit:battery}{电池}时的正确连接方式
        ——并说明为何无需先搭建电路。
  \item 抽屉中刚取出的单体\omterm{def:g3:first-electric-circuit:battery}{电池}读数如下：
        \qty{1.58}{V}、\qty{1.32}{V}、\qty{0.9}{V}、
        \qty{0.1}{V}。请将其分类：新鲜、可用、电量偏低、已耗尽。
  \item 一节九伏\omterm{def:g3:first-electric-circuit:battery}{电池}的读数为 \qty{8.9}{V}。大致
        每节内部单体\omterm{def:g3:first-electric-circuit:battery}{电池}贡献多少？
  \item 为何一节将耗尽的单体\omterm{def:g3:first-electric-circuit:battery}{电池}仍显示\emph{一定}
        \omterm{def:g8:voltage-voltmeter:voltage}{电压}，尽管它已无法点亮任何东西？（它的
        推力仍微弱地存在；它失去了何种能力，无法
        驱动？）
\end{enumerate}

\medskip\noindent\textbf{第二部分——小物件架。}
\begin{enumerate}[resume]
  \item 相机需要两节\omterm{def:g3:first-electric-circuit:battery}{电池}首尾相接，它需要多大\omterm{def:g8:voltage-voltmeter:voltage}{电压}，
        以及第 2 题中剩下的那两节
        \qty{1.32}{V} \omterm{def:g3:first-electric-circuit:battery}{电池}能提供多少？
  \item 旧收音机要用四节\omterm{def:g3:first-electric-circuit:battery}{电池}首尾相接，
        \omterm{def:g8:voltage-voltmeter:voltage}{电压}低于 \qty{5}{V} 就罢工。抽屉里的\omterm{def:g3:first-electric-circuit:battery}{电池}中，
        哪些（如果有的话）能让它工作？
  \item 弟弟把三节\omterm{def:g3:first-electric-circuit:battery}{电池}硬塞进只能装两节的
        手电筒里，说是“为了更亮”，还把第三节横卡在
        弹簧上。\omterm{def:g3:first-electric-circuit:bulb}{灯泡}\omterm{def:g8:voltage-voltmeter:nominal}{额定电压}为 \qty{3}{V}。
        请用\omterm{def:g8:voltage-voltmeter:nominal}{额定电压}的语言预测会出现什么情况。
  \item 门铃变压器的铭牌上写着“\qty{8}{V}”。
        有人提议把门铃直接接到一节九伏\omterm{def:g3:first-electric-circuit:battery}{电池}上，
        理由是“八和九差不多”。
        请评估：这是合理搭配，还是在慢慢损坏门铃？
\end{enumerate}

\medskip\noindent\textbf{第三部分——审计定律。}
\begin{enumerate}[resume]
  \item 在手电筒中（\omterm{def:g3:first-electric-circuit:bulb}{灯泡}与\omterm{def:g3:first-electric-circuit:battery}{电池}\omterm{def:g4:series-circuits:series}{串联}，\omterm{ex:g3:first-electric-circuit:switch}{开关}断开），
        审计员测得：断开的\omterm{ex:g3:first-electric-circuit:switch}{开关}两端
        \omterm{def:g8:voltage-voltmeter:voltage}{电压}为 \qty{3.1}{V}；\omterm{def:g3:first-electric-circuit:bulb}{灯泡}两端\omterm{def:g8:voltage-voltmeter:voltage}{电压}为 \qty{0}{V}。请解释两个读数，
        并预测\omterm{ex:g3:first-electric-circuit:switch}{开关}闭合（咔嗒一声）后两者各为多少。
  \item 在手电筒线路的每一段两端，
        电表读数均为 \qty{0}{V}，无论灯亮还是灯灭。
        用“推力”的语言说明，这对良好导线意味着什么？
  \item 按家庭的既定规矩，市电插座\emph{不}在审计之列。
        请用本章给出的数值重新表述这一定律，
        并写出其背后的算式。
  \item 写审计报告，共三句——\omterm{def:g8:voltage-voltmeter:voltmeter}{电压表}是什么、又绝不是什么；
        标签上\omterm{def:g8:voltage-voltmeter:nominal}{额定电压}承诺了什么、又威胁了什么；
        以及在任何闲置电路中，\omterm{def:g8:voltage-voltmeter:voltage}{电压}在何处等待。
        ——
\end{enumerate}
\end{problem}
